\documentclass[12pt, a4paper]{article}
\usepackage[utf8]{inputenc}
\usepackage[T1]{fontenc}
\usepackage[french]{babel}
\usepackage{amsmath, amssymb, amsthm}
\usepackage{geometry}
\usepackage{listings}
\geometry{margin=2.5cm}

\newtheorem{theorem}{Théorème}[section]
\newtheorem{lemma}[theorem]{Lemme}
\newtheorem{definition}[theorem]{Définition}

\title{Sur la Conjecture d'Erdős-Straus : Analyse Algébrique et Décomposition Modulaire}
\author{Charles EDOU NZE\thanks{Charles EDOU NZE, chercheur indépendant}}
\date{\today}

\lstdefinelanguage{lean}{
  keywords={import, def, theorem, lemma, by, sorry, Prop, Nat, open, section, Exists, fun},
  sensitive=true,
  comment=[l]--
}

\begin{document}
\maketitle

\begin{abstract}
Ce document présente une analyse rigoureuse et une décomposition structurelle de la conjecture d'Erdős-Straus. Nous y exposons les définitions axiomatiques, passons en revue la littérature récente pertinente, isolons des lemmes clés pour des classes de congruences spécifiques, et fournissons une architecture d'autoformalisation pour un assistant de preuve tel que Lean 4.
\end{abstract}

\tableofcontents
\newpage

\section{Définitions Axiomatiques et Contexte}
\begin{definition}
Pour tout entier $n \in \mathbb{N}$ avec $n \ge 2$, l'équation d'Erdős-Straus est définie comme l'équation diophantienne :
\[\frac{4}{n} = \frac{1}{x} + \frac{1}{y} + \frac{1}{z}\]
où $x, y, z \in \mathbb{Z}_{>0}$.
\end{definition}

\subsection{Recherche de Littérature Contextuelle}
Le problème, formulé par Paul Erdős et Ernst G. Straus en 1948, postule que l'équation admet toujours une solution entière positive pour $n \ge 2$. Des études récentes par des auteurs tels que Miguel Angel Lopez ont classifié les solutions en types, comme le Type A et le Type B, en définissant un système complet de congruences. De plus, Philemon Urbain Mballa a exploré une connexion inattendue entre la fonction zêta discrète et la conjecture d'Erdős-Straus via une décomposition additive. Cette méthodologie présente de fortes analogies avec les manipulations algébriques utilisées pour résoudre les équations de Pell-Fermat, où les solutions sont regroupées par invariants modulaires et analysées par des descentes infinies ou des classifications structurelles.

\section{Stratégie de Preuve et Isolation des Lemmes}
L'approche fondamentale consiste à partitionner les entiers $n$ par leur classe de congruence modulo un entier, historiquement $840$, bien que des cas généraux puissent être étudiés pour les nombres premiers.

\begin{lemma}
Pour $n = 4k+3$ avec $k \in \mathbb{N}$, l'équation admet une solution entière positive.
\end{lemma}
\begin{proof}
Soit $n = 4k+3$.
Nous posons $x = (n+1)/4 = (4k+4)/4 = k+1$. Comme $n \ge 3$, $k \ge 0$, et donc $x \ge 1$, ce qui garantit que $x \in \mathbb{Z}_{>0}$.
Nous évaluons la différence restante :
\begin{align*}
\frac{4}{n} - \frac{1}{x} &= \frac{4}{n} - \frac{1}{k+1} \\
&= \frac{4(k+1) - n}{n(k+1)} \\
&= \frac{4k+4 - (4k+3)}{n(k+1)} \\
&= \frac{1}{n(k+1)}
\end{align*}
La fraction restante $\frac{1}{n(k+1)}$ doit être divisée en deux fractions unitaires $\frac{1}{y} + \frac{1}{z}$.
Nous utilisons l'identité standard de séparation des fractions égyptiennes :
\[ \frac{1}{A} = \frac{1}{A+1} + \frac{1}{A(A+1)} \]
En appliquant ceci à $A = n(k+1)$, nous posons :
$y = n(k+1) + 1$
et
$z = n(k+1)(n(k+1)+1)$
Puisque $n \ge 3$ et $k \ge 0$, $y$ et $z$ sont tous deux des entiers strictement positifs.
Par substitution :
\begin{align*}
\frac{1}{x} + \frac{1}{y} + \frac{1}{z} &= \frac{1}{k+1} + \frac{1}{n(k+1) + 1} + \frac{1}{n(k+1)(n(k+1)+1)} \\
&= \frac{1}{k+1} + \frac{n(k+1)+1 + 1}{(n(k+1)+1)(n(k+1)(n(k+1)+1))} \\
\end{align*}
La somme directe des deux derniers termes est exactement $\frac{1}{n(k+1)}$. Démontrons-le explicitement :
\begin{align*}
\frac{1}{y} + \frac{1}{z} &= \frac{1}{n(k+1)+1} + \frac{1}{n(k+1)(n(k+1)+1)} \\
&= \frac{n(k+1)}{n(k+1)(n(k+1)+1)} + \frac{1}{n(k+1)(n(k+1)+1)} \\
&= \frac{n(k+1)+1}{n(k+1)(n(k+1)+1)} \\
&= \frac{1}{n(k+1)}
\end{align*}
En ajoutant $\frac{1}{x} = \frac{1}{k+1}$, nous obtenons :
\begin{align*}
\frac{1}{x} + \frac{1}{y} + \frac{1}{z} &= \frac{1}{k+1} + \frac{1}{n(k+1)} \\
&= \frac{n}{n(k+1)} + \frac{1}{n(k+1)} \\
&= \frac{n+1}{n(k+1)} \\
&= \frac{4k+4}{n(k+1)} \\
&= \frac{4(k+1)}{n(k+1)} \\
&= \frac{4}{n}
\end{align*}
Ceci prouve le lemme pour tout $n \equiv 3 \pmod 4$.
\end{proof}

\section{Architecture pour l'Autoformalisation}
La formalisation de la conjecture d'Erdős-Straus nécessite de structurer l'énoncé et la décomposition de l'espace de recherche en Lean 4.

\begin{lstlisting}[language=lean, basicstyle=\ttfamily\small, breaklines=true]
import Mathlib.Data.Nat.Basic
import Mathlib.Tactic.Omega
import Mathlib.Tactic.Ring

-- Definition axiomatique de la propriete d'Erdos-Straus
def SatisfiesErdosStraus (n : Nat) : Prop :=
  Exists (fun x => Exists (fun y => Exists (fun z => x > 0 /\ y > 0 /\ z > 0 /\ 4 * x * y * z = n * (y * z + x * z + x * y))))

-- Lemme pour mod 4 = 3
lemma erdos_straus_mod_4_3 (k : Nat) : SatisfiesErdosStraus (4 * k + 3) := by
  sorry

-- Theoreme principal
theorem erdos_straus_conjecture (n : Nat) (hn : n >= 2) : SatisfiesErdosStraus n := by
  sorry
\end{lstlisting}

\end{document}
